
\section{证明纲要，\trFloerHomology}
\label{sec:outline-floer-homology}

我们将用来证明\trArnoldConjecture（在一个限制性条件下）的主要工具，是\trFloerHomology。\trFloer{}的想法 [28, 29, 30, 31]，是要去“计数”$H$ 的诸\trPeriodic{}解，也就是它的\trFlow{}在时刻 $t=1$ 的诸\trFixedPoint，从而以一种类似于曾使我们得以给出\trManifold{}上一个\trMorseFunction{}的\trCriticalPoint{}数目之下界的方式，给出它们个数的一个下界（\cref{thm:morse-inequalities}）。

\begin{enumerate}[label=(\arabic*)]
\item 我们首先在一个适当的空间上（这个空间可以看成是一个无穷维\trManifold；见\cref{subsec:loop-space}）定义一个“函数”，即\textbf{\trActionFunctional}（\cref{sec:action-functional}）。这个泛函的\trCriticalPoint{}将恰好就是所欲求的那些\trPeriodic{}解。
\item 借助这些\trCriticalPoint，我们将构造一个\trComplex。我们先取由它们生成的 $\ZZ/2$ 上的诸\trVectorSpace。为了在这些空间上定义一个\trGraded，我们将需要\trCriticalPoint{}的\trIndex{}的一个类似物，我们将在第 7 章中定义它；而当\trHamiltonian{}量不依赖时间时，它就是\trMorse{}意义下的那个\trIndex（至多相差一个平移）（推论~7.2.2）。
\item 为了定义这个\trComplex{}的\trDifferential，我们将使用一个\trVectorField，即该泛函的\trGradient（或其负）；我们将在\cref{sec:gradient-floer-equation}中描述它，连同定义它的那个\trMetric{}一起。
\item 这时的想法便是去“计数”它的诸\trTrajectory。在\cref{subsec:definition-energy}中，我们将凸显一个性质——\trEnergy{}的有限性；对于\trCompact\trManifold{}上的一个\trPseudogradient{}场而言，它是自动满足的，而在这里它将派上用场。
\item 我们将证明，有限\trEnergy{}的诸\trTrajectory{}的确连接两个\trCriticalPoint（\cref{thm:trajectories-tend-critical-points}）。
\item 我们将需要有限\trEnergy{}的诸\trTrajectory{}所成空间的一个\trCompact{}性性质（\trGromov{}\trCompact{}性，\cref{thm:gromov-compactness-theorem}）。
\item 正如在\trMorse{}情形中那样（\cref{subsec:compactness}），我们将需要另一个\trCompact{}性性质来定义\trDifferential{} $\partial$（在 §9.1.c 中）。
\item 当然，我们还需要这些\trTrajectory{}所成的空间都是\trManifold，也需要一个\trGenericity{}性质，诸如对充分接近的\trVectorField{}而言的\trSmale{}性质（\cref{thm:smale-theorem}）；我们将在第 8 章中做到这一点。
\item 这样定义了所设想的\trComplex{}的\trDifferential{}之后，我们仍需验证 $\partial\circ\partial=0$。在\trMorse{}情形中，我们用的是一个\trGluing{}性质（\cref{subsec:manifold-structure-broken}）。我们将在第 9 章中证明一个类似的性质。
\item 利用与\cref{sec:homology-independence}中相同的方法，我们将证明，我们最终定义的这个\trComplex{}的\trHomology{}既不依赖于泛函的选取，也不依赖于\trVectorField{}的选取。
\item 最后，通过考察\trAutonomous{}的 $\cC^2$-小\trHamiltonian{}量的情形，这个不依赖性将告诉我们：我们由此构造出的那个\trHomology，即\trFloerHomology，与该\trManifold{}的\trMorseHomology{}相重合。参与这个\trComplex{}的诸\trVectorSpace{}的维数之间那个显见的不等式，将以一种完全类似于\trMorse{}不等式的方式，给出所期待的结果（\cref{thm:morse-inequalities}）。
\end{enumerate}

\textbf{诸假设}

我们将假设\trSymplectic\trManifold{} $W$ 满足下面两个条件。

\begin{assumption}\label{thm:symplectic-asphericity}
对每个 $\cC^\infty$ 映射 $w\colon S^2\to W$，都有
\[
\int_{S^2}w^\star\omega=0\text{。}
\]
\end{assumption}

在文献中，这个条件有时被称为“\textbf{\trSymplectic\trAsphericity}”（symplectic asphericity，即\trSymplecticForm{}在诸\trSphere{}上为零）。它可以表述为
\[
\langle\omega,\pi_2(W)\rangle=0\text{。}
\]

\begin{assumption}\label{thm:symplectic-trivialization}
对每个 $\cC^\infty$ 映射 $w\colon S^2\to W$，\trFiberBundle{} $w^\star TW$ 都存在一个\textbf{\trSymplectic{}\trTrivialization}（symplectic trivialization）。
\end{assumption}

这个条件可以表述为
\[
\langle c_1(TW),\pi_2(W)\rangle=0\text{，}
\]
其中 $c_1(TW)$ 表示复\trVectorBundle{} $TW$ 的\trFirstChernClass（参见\cref{sec:appB-chern-classes}）。

这两个条件在下述情形下都是满足的，例如：当从 $S^2$ 到 $W$ 的每个映射都能扩张到球体 $B^3$ 上时，也就是说，当 $\pi_2(W)=0$ 时。
